\documentclass[12pt]{article}

% Packages
\usepackage[margin=1in]{geometry}
\usepackage{amsmath}
\usepackage{amssymb}
\usepackage{hyperref}

% Document metadata
\title{A Simple \LaTeX{} Document}
\author{Your Name}
\date{\today}

\begin{document}

\maketitle

\tableofcontents
\newpage

% -----------------------------------------------
\section{Introduction}
% -----------------------------------------------

Welcome to \LaTeX{}! This is a sample document demonstrating common features.
\LaTeX{} is a powerful typesetting system widely used in academia and publishing.

% -----------------------------------------------
\section{Text Formatting}
% -----------------------------------------------

You can make text \textbf{bold}, \textit{italic}, or \underline{underlined}.
You can also combine them: \textbf{\textit{bold italic}}.

Here is a new paragraph. In \LaTeX{}, paragraphs are separated by a blank line.

% -----------------------------------------------
\section{Lists}
% -----------------------------------------------

\subsection{Unordered List}
\begin{itemize}
  \item First item
  \item Second item
  \item Third item
\end{itemize}

\subsection{Ordered List}
\begin{enumerate}
  \item Step one
  \item Step two
  \item Step three
\end{enumerate}

% -----------------------------------------------
\section{Mathematics}
% -----------------------------------------------

Inline math: The quadratic formula is $x = \frac{-b \pm \sqrt{b^2 - 4ac}}{2a}$.

Display math:
\[
  \int_{-\infty}^{\infty} e^{-x^2} \, dx = \sqrt{\pi}
\]

An aligned equation:
\begin{align}
  E &= mc^2 \\
  F &= ma
\end{align}

% -----------------------------------------------
\section{Tables}
% -----------------------------------------------

\begin{table}[h]
  \centering
  \begin{tabular}{|l|c|r|}
    \hline
    \textbf{Name} & \textbf{Age} & \textbf{Score} \\
    \hline
    Alice   & 24 & 95 \\
    Bob     & 30 & 87 \\
    Charlie & 22 & 91 \\
    \hline
  \end{tabular}
  \caption{Sample data table}
  \label{tab:sample}
\end{table}

% -----------------------------------------------
\section{Conclusion}
% -----------------------------------------------

This covers the basics of \LaTeX{}. From here you can explore packages for
figures, bibliography management with BibTeX, custom macros, and much more.

\end{document}